% GENERATED FILE. Edit canonical lessons, then run npm run generate:books.
% canonical-source-hash: fnv1a64:90856cf278272947
% canonical-lessons: JA-C88-utsusu, JA-C88-tamesu, JA-C88-modosu, JA-C88-nokosu, JA-C88-taosu

\chapter{To copy, To try out, To put back, To leave behind, To knock down}
\label{ch:ja-to-copy-to-try-out-to-put-back-to-leave-behind-to-knock-down}
\hlchaptermodality{88}

\begin{chapteropening}
\textbf{By the end of this chapter:} \emph{I can say to copy, to try out, to put back, to leave behind and to knock down.}

Complete the last lesson of chapter 88: I can say to copy, to try out, to put back, to leave behind and to knock down.
\end{chapteropening}

\section[\texorpdfstring{\ja{うつす} (utsusu)}{うつす (utsusu)}]{\ja{うつす} (utsusu) — to copy}

\label{lesson:JA-C88-utsusu}



\begin{quote}
Before the new one: say the Japanese for to hang out to dry, then the Japanese for to point; to hold up.
\end{quote}

\subsection*{You'll want to know: \ja{うつす}}
\textbf{\ja{うつす}} — \emph{utsusu} — \textquotedblleft{}to copy\textquotedblright{}.

\begin{grammarlens}[title={What you've built}]
One more word for this chapter.
\end{grammarlens}

\subsection*{Your turn}
\begin{itemize}
\raggedright
  \item \emph{Say it:} \emph{utsusu}
  \item \emph{Say it:} \emph{utsusu}, once more
  \item \emph{Say it:} say \emph{utsusu}
  \item \emph{Recall:} say the Japanese for to hang out to dry, then the Japanese for to point; to hold up, then say \emph{utsusu} again
\end{itemize}

\begin{tcolorbox}[breakable,colback=teal!4,colframe=teal!35!black,title={Before you move on}]
What does \ja{うつす} mean? (\textquotedblleft{}To copy\textquotedblright{}.) Say it once more.
\end{tcolorbox}
\section[\texorpdfstring{\ja{ためす} (tamesu)}{ためす (tamesu)}]{\ja{ためす} (tamesu) — to try out}

\label{lesson:JA-C88-tamesu}



\begin{quote}
Before the new one: say the Japanese for to point; to hold up, then the Japanese for to copy.
\end{quote}

\subsection*{You'll want to know: \ja{ためす}}
\textbf{\ja{ためす}} — \emph{tamesu} — \textquotedblleft{}to try out\textquotedblright{}.

\begin{grammarlens}[title={What you've built}]
One more word for this chapter.
\end{grammarlens}

\subsection*{Your turn}
\begin{itemize}
\raggedright
  \item \emph{Say it:} \emph{tamesu}
  \item \emph{Say it:} \emph{tamesu}, once more
  \item \emph{Say it:} say \emph{tamesu}
  \item \emph{Recall:} say the Japanese for to point; to hold up, then the Japanese for to copy, then say \emph{tamesu} again
\end{itemize}

\begin{tcolorbox}[breakable,colback=teal!4,colframe=teal!35!black,title={Before you move on}]
What does \ja{ためす} mean? (\textquotedblleft{}To try out\textquotedblright{}.) Say it once more.
\end{tcolorbox}
\section[\texorpdfstring{\ja{もどす} (modosu)}{もどす (modosu)}]{\ja{もどす} (modosu) — to put back}

\label{lesson:JA-C88-modosu}



\begin{quote}
Before the new one: say the Japanese for to copy, then the Japanese for to try out.
\end{quote}

\subsection*{You'll want to know: \ja{もどす}}
\textbf{\ja{もどす}} — \emph{modosu} — \textquotedblleft{}to put back\textquotedblright{}.

\begin{grammarlens}[title={What you've built}]
One more word for this chapter.
\end{grammarlens}

\subsection*{Your turn}
\begin{itemize}
\raggedright
  \item \emph{Say it:} \emph{modosu}
  \item \emph{Say it:} \emph{modosu}, once more
  \item \emph{Say it:} say \emph{modosu}
  \item \emph{Recall:} say the Japanese for to copy, then the Japanese for to try out, then say \emph{modosu} again
\end{itemize}

\begin{tcolorbox}[breakable,colback=teal!4,colframe=teal!35!black,title={Before you move on}]
What does \ja{もどす} mean? (\textquotedblleft{}To put back\textquotedblright{}.) Say it once more.
\end{tcolorbox}
\section[\texorpdfstring{\ja{のこす} (nokosu)}{のこす (nokosu)}]{\ja{のこす} (nokosu) — to leave behind}

\label{lesson:JA-C88-nokosu}



\begin{quote}
Before the new one: say the Japanese for to try out, then the Japanese for to put back.
\end{quote}

\subsection*{You'll want to know: \ja{のこす}}
\textbf{\ja{のこす}} — \emph{nokosu} — \textquotedblleft{}to leave behind\textquotedblright{}.

\begin{grammarlens}[title={What you've built}]
One more word for this chapter.
\end{grammarlens}

\subsection*{Your turn}
\begin{itemize}
\raggedright
  \item \emph{Say it:} \emph{nokosu}
  \item \emph{Say it:} \emph{nokosu}, once more
  \item \emph{Say it:} say \emph{nokosu}
  \item \emph{Recall:} say the Japanese for to try out, then the Japanese for to put back, then say \emph{nokosu} again
\end{itemize}

\begin{tcolorbox}[breakable,colback=teal!4,colframe=teal!35!black,title={Before you move on}]
What does \ja{のこす} mean? (\textquotedblleft{}To leave behind\textquotedblright{}.) Say it once more.
\end{tcolorbox}
\section[\texorpdfstring{\ja{たおす} (taosu)}{たおす (taosu)}]{\ja{たおす} (taosu) — to knock down}

\label{lesson:JA-C88-taosu}



\begin{quote}
Before the new one: say the Japanese for to put back, then the Japanese for to leave behind.
\end{quote}

\subsection*{You'll want to know: \ja{たおす}}
\textbf{\ja{たおす}} — \emph{taosu} — \textquotedblleft{}to knock down\textquotedblright{}.

\begin{grammarlens}[title={What you've built}]
One more word for this chapter.
\end{grammarlens}

\subsection*{Your turn}
\begin{itemize}
\raggedright
  \item \emph{Say it:} \emph{taosu}
  \item \emph{Say it:} \emph{taosu}, once more
  \item \emph{Say it:} say \emph{taosu}
  \item \emph{Recall:} say the Japanese for to put back, then the Japanese for to leave behind, then say \emph{taosu} again
\end{itemize}

\begin{tcolorbox}[breakable,colback=teal!4,colframe=teal!35!black,title={Before you move on}]
What does \ja{たおす} mean? (\textquotedblleft{}To knock down\textquotedblright{}.) Say it once more.
\end{tcolorbox}
